\subsection{Impedanz der Spule}
\label{subsec:impedanz-der-spule}

Neben dem Kondensator besitzt auch die Spule einen frequenzabhängigen
Wechselstromwiderstand. Dieser wird als induktiver Blindwiderstand bezeichnet.
Während der Kondensator Stromänderungen über seine elektrische Ladung beschreibt,
beruht das Verhalten der Spule auf dem Aufbau und Abbau ihres Magnetfeldes.

Die Grundgleichung der idealen Spule lautet

\begin{equation}
	u_L(t) = L \frac{di(t)}{dt}.
	\label{eq:spule-grundgleichung}
\end{equation}

Dabei ist \(u_L(t)\) die Spannung an der Spule, \(i(t)\) der Strom durch die
Spule und \(L\) die Induktivität.

\begin{PhysicsBox}[Grundidee der Spule]
	\label{box:grundidee-der-spule}
	Eine Spule widersetzt sich Änderungen des Stroms. Ändert sich der Strom,
	ändert sich auch das Magnetfeld der Spule. Dieses veränderliche Magnetfeld
	erzeugt eine Gegenspannung. Deshalb reagiert die Spule besonders stark auf
	schnelle Stromänderungen.
\end{PhysicsBox}

\subsubsection{Herleitung mit sinusförmigem Strom}
\label{subsubsec:spule-sinusfoermiger-strom}

Wir betrachten einen sinusförmigen Strom

\begin{equation}
	i(t) = \hat{I} \sin(\omega t),
	\label{eq:spule-strom-sinus}
\end{equation}

wobei \(\hat{I}\) der Scheitelwert des Stroms und

\begin{equation}
	\omega = 2\pi f
	\label{eq:kreisfrequenz-spule}
\end{equation}

die Kreisfrequenz ist.

Die Spannung an der Spule ergibt sich aus der Ableitung des Stroms:

\begin{align}
	u_L(t)
	&= L \frac{di(t)}{dt} \\[4pt]
	&= L \frac{d}{dt} \left( \hat{I} \sin(\omega t) \right) \\[4pt]
	&= L \hat{I} \omega \cos(\omega t).
	\label{eq:spule-spannung-cosinus}
\end{align}

Damit folgt

\begin{equation}
	u_L(t) = \omega L \hat{I} \cos(\omega t).
	\label{eq:spule-spannung-omega-l}
\end{equation}

Da gilt

\begin{equation}
	\cos(\omega t)
	=
	\sin\left(\omega t + \frac{\pi}{2}\right),
	\label{eq:cosinus-als-verschobener-sinus}
\end{equation}

kann man schreiben:

\begin{equation}
	u_L(t)
	=
	\omega L \hat{I}
	\sin\left(\omega t + \frac{\pi}{2}\right).
	\label{eq:spule-spannung-phasenverschoben}
\end{equation}

Die Spannung eilt dem Strom also um \(90^\circ\) voraus.

\begin{DidacticBox}[Merksatz zur Spule]
	\label{box:merksatz-spule-strom-hinkt-nach}
	Bei der Spule eilt die Spannung dem Strom um \(90^\circ\) voraus.
	
	Anders gesagt: Der Strom hinkt der Spannung um \(90^\circ\) nach.
\end{DidacticBox}

\subsubsection{Induktiver Blindwiderstand}
\label{subsubsec:induktiver-blindwiderstand}

Aus Gleichung~\eqref{eq:spule-spannung-omega-l} liest man den Scheitelwert
der Spannung ab:

\begin{equation}
	\hat{U} = \omega L \hat{I}.
	\label{eq:spule-scheitelwert-spannung}
\end{equation}

Daraus folgt für das Verhältnis von Spannung zu Strom:

\begin{equation}
	\frac{\hat{U}}{\hat{I}} = \omega L.
	\label{eq:spule-verhaeltnis-u-i}
\end{equation}

Der Betrag des induktiven Blindwiderstands lautet daher

\begin{equation}
	X_L = \omega L.
	\label{eq:induktiver-blindwiderstand}
\end{equation}

Mit \(\omega = 2\pi f\) erhält man

\begin{equation}
	X_L = 2\pi f L.
	\label{eq:induktiver-blindwiderstand-frequenz}
\end{equation}

\begin{MathBox}[Induktiver Blindwiderstand]
	\label{box:induktiver-blindwiderstand}
	Der induktive Blindwiderstand einer idealen Spule ist
	
	\[
	X_L = \omega L = 2\pi f L.
	\]
	
	Er wächst proportional mit der Frequenz \(f\) und mit der Induktivität \(L\).
\end{MathBox}

Die Einheit ist Ohm:

\begin{equation}
	[X_L]
	=
	[\omega L]
	=
	\frac{1}{\mathrm{s}} \cdot \mathrm{H}
	=
	\Omega.
	\label{eq:einheit-induktiver-blindwiderstand}
\end{equation}

\subsubsection{Komplexe Herleitung der Impedanz}
\label{subsubsec:komplexe-herleitung-impedanz-spule}

In der komplexen Wechselstromrechnung beschreibt man sinusförmige Größen mit
dem komplexen Exponentialansatz. Für den Strom schreiben wir

\begin{equation}
	i(t) = \hat{I} e^{j\omega t}.
	\label{eq:spule-komplexer-strom}
\end{equation}

Die Spulenspannung ergibt sich wieder aus

\begin{equation}
	u_L(t) = L \frac{di(t)}{dt}.
	\label{eq:spule-komplexe-grundgleichung}
\end{equation}

Einsetzen liefert

\begin{align}
	u_L(t)
	&= L \frac{d}{dt}
	\left(
	\hat{I} e^{j\omega t}
	\right) \\[4pt]
	&= L \cdot j\omega \hat{I} e^{j\omega t} \\[4pt]
	&= j\omega L \hat{I} e^{j\omega t}.
	\label{eq:spule-komplexe-spannung}
\end{align}

Da

\begin{equation}
	i(t) = \hat{I} e^{j\omega t}
	\label{eq:spule-komplexer-strom-wiederholung}
\end{equation}

gilt, folgt unmittelbar:

\begin{equation}
	u_L(t) = j\omega L \, i(t).
	\label{eq:spule-u-gleich-z-mal-i}
\end{equation}

Die Impedanz ist definiert als Verhältnis von Spannung zu Strom:

\begin{equation}
	Z_L = \frac{U}{I}.
	\label{eq:definition-impedanz-spule}
\end{equation}

Damit erhält man für die ideale Spule

\begin{equation}
	Z_L = j\omega L.
	\label{eq:impedanz-spule}
\end{equation}

Oder mit \(\omega = 2\pi f\):

\begin{equation}
	Z_L = j 2\pi f L.
	\label{eq:impedanz-spule-frequenz}
\end{equation}

\begin{MathBox}[Impedanz der idealen Spule]
	\label{box:impedanz-ideale-spule}
	Die komplexe Impedanz einer idealen Spule lautet
	
	\[
	Z_L = j\omega L.
	\]
	
	Der Faktor \(j\) beschreibt die Phasenverschiebung um \(+90^\circ\).
	Die Spannung eilt dem Strom voraus.
\end{MathBox}

\subsubsection{Betrag und Phase}
\label{subsubsec:betrag-und-phase-spule}

Die Impedanz der Spule lautet

\begin{equation}
	Z_L = j\omega L.
	\label{eq:spule-impedanz-betrag-phase-ausgangspunkt}
\end{equation}

Der Betrag ist

\begin{equation}
	|Z_L| = \omega L.
	\label{eq:betrag-impedanz-spule}
\end{equation}

Die Phase beträgt

\begin{equation}
	\varphi = +90^\circ.
	\label{eq:phase-impedanz-spule}
\end{equation}

Damit kann man die Impedanz auch in Polarform schreiben:

\begin{equation}
	Z_L = \omega L \, e^{j\frac{\pi}{2}}.
	\label{eq:spule-impedanz-polarform}
\end{equation}

\subsubsection{Vergleich mit Widerstand und Kondensator}
\label{subsubsec:vergleich-r-l-c-impedanzen}

Zum Vergleich:

\begin{align}
	Z_R &= R,
	\label{eq:impedanz-widerstand-vergleich} \\[4pt]
	Z_L &= j\omega L,
	\label{eq:impedanz-spule-vergleich} \\[4pt]
	Z_C &= \frac{1}{j\omega C}.
	\label{eq:impedanz-kondensator-vergleich}
\end{align}

Der Widerstand besitzt keine Phasenverschiebung. Die Spule besitzt eine
positive Phasenverschiebung von \(+90^\circ\), der Kondensator eine negative
Phasenverschiebung von \(-90^\circ\).

\begin{center}
	\begin{tabular}{c|c|c|c}
		Bauteil & Grundgleichung & Impedanz & Phase \\ \hline
		Widerstand & \(u = Ri\) & \(Z_R = R\) & \(0^\circ\) \\
		Spule & \(u_L = L\dfrac{di}{dt}\) & \(Z_L = j\omega L\) & \(+90^\circ\) \\
		Kondensator & \(i_C = C\dfrac{du}{dt}\) & \(Z_C = \dfrac{1}{j\omega C}\) & \(-90^\circ\)
	\end{tabular}
\end{center}

\begin{DidacticBox}[Spule und Kondensator im Vergleich]
	\label{box:spule-kondensator-vergleich}
	Bei der Spule wächst der Blindwiderstand mit der Frequenz:
	
	\[
	X_L = \omega L.
	\]
	
	Beim Kondensator sinkt der Blindwiderstand mit der Frequenz:
	
	\[
	X_C = \frac{1}{\omega C}.
	\]
	
	Die Spule sperrt hohe Frequenzen stärker, der Kondensator lässt hohe
	Frequenzen leichter passieren.
\end{DidacticBox}

\subsubsection{Grenzfälle}
\label{subsubsec:grenzfaelle-spule}

Für Gleichstrom gilt

\begin{equation}
	f = 0
	\label{eq:gleichstromfrequenz-spule}
\end{equation}

und damit

\begin{equation}
	X_L = 2\pi f L = 0.
	\label{eq:spule-gleichstrom-blindwiderstand}
\end{equation}

Eine ideale Spule besitzt bei Gleichstrom also keinen induktiven
Blindwiderstand. In der Realität bleibt jedoch der ohmsche Widerstand des
Kupferdrahtes erhalten.

Für sehr hohe Frequenzen wird der Blindwiderstand dagegen groß:

\begin{equation}
	X_L = 2\pi f L \rightarrow \infty
	\quad \text{für} \quad f \rightarrow \infty.
	\label{eq:spule-hohe-frequenzen}
\end{equation}

\begin{PhysicsBox}[Physikalische Bedeutung]
	\label{box:physikalische-bedeutung-spulenimpedanz}
	Eine ideale Spule lässt Gleichstrom nach dem Einschwingvorgang nahezu
	ungehindert fließen. Wechselstrom wird dagegen umso stärker behindert, je
	höher seine Frequenz ist. Ursache ist die Selbstinduktion der Spule.
\end{PhysicsBox}

\subsubsection{Ergebnis}
\label{subsubsec:ergebnis-impedanz-spule}

Die Impedanz der idealen Spule lautet

\begin{equation}
	\boxed{
		Z_L = j\omega L
	}
	\label{eq:endformel-impedanz-spule}
\end{equation}

mit

\begin{equation}
	\omega = 2\pi f.
	\label{eq:endformel-kreisfrequenz-spule}
\end{equation}

Damit ergibt sich

\begin{equation}
	\boxed{
		Z_L = j2\pi f L
	}
	\label{eq:endformel-impedanz-spule-frequenz}
\end{equation}

und für den Betrag

\begin{equation}
	\boxed{
		|Z_L| = X_L = 2\pi f L.
	}
	\label{eq:endformel-betrag-impedanz-spule}
\end{equation}

\begin{DidacticBox}[Kurzer Merksatz]
	\label{box:kurzer-merksatz-spule}
	Spule:
	
	\[
	Z_L = j\omega L.
	\]
	
	Die Spannung eilt vor, der Strom hinkt nach.
\end{DidacticBox}